\section{Conclusions}
\label{sec:conclusions}

\subsection{Summary of findings}

This study screens European residential heat across 29 markets to 2050 in three operating roles. The tests address distinct questions but share some input assumptions. Seasonal-storage assumptions affect the base-load and dispatch rankings. Conditional on entry, however, the capital-recovery ratio contains no fuel- or storage-cost term; it depends on utilisation, scarcity margin, WACC, fixed O\&M and turbine CAPEX. Four findings follow.

(1) On central 2050 levelised cost, the best heat pump delivers at \euro{}118.5/MWh,
compared with \euro{}122.8/MWh for the hydrogen boiler and \euro{}132.7/MWh for gas. Heat
pumps beat gas in 24 of the 29 markets under the central carbon path; the count ranges
from 19 to 29 across the carbon paths tested. Against hydrogen, heat pumps win in 22
markets when firm-generation cost is allocated pro rata to heating's annual
electricity-demand share. Allocation by the change in peaking capacity gives 19 heat-pump
wins at the low bound and 17 at the high bound. The median 2050 country gap is
\euro{}15.4/MWh, while the mean weighted by 2025 country heat demand is \euro{}1.4/MWh
because the hydrogen-favouring markets are relatively large.

(2) Under the central symmetric accounting, hydrogen wins the base-load comparison in
seven markets and heat pumps in 22. Across the hydrogen-price paths tested, the number of
heat-pump wins ranges from 13 to 29. Dickel \citep{Dickel2024ET29} concludes that the
hydrogen realistically supplied to heat would be blue, at about \euro{}81/MWh on the route
priced here \citep{IEAGHG2022blue}. A uniform blue-hydrogen price of \euro{}81/MWh removes
hydrogen's advantage in all seven central green-hydrogen base-load markets. Scaling each
country's green-hydrogen price by the 81/50 ratio instead leaves Denmark as the sole
hydrogen win. Levelling the carriers' fiscal treatment moves the comparison in the same
direction, by an amount outside this study's scope.

(3) Repriced on SRMC, on the high carbon path and the equal-efficiency
series, hydrogen wins the power peak in a different seven. All seven are salt-cavern markets on a binary geology screen, and alternative definitions give six or nine, so the count is a screening result. The
peaking increment is about 278~GW, roughly 31~per~cent of the study area's 889~GW peak, yet supplies about
0.86~per~cent of its electricity. At a hydrogen-ready turbine capital cost of \euro{}600/kW, no tested market reaches full
peaker-capital recovery across the assumed scarcity margins. Capacity-weighted recovery is
about 9~per~cent at the central 2050 margin and ranges from 4.4 to 17.4~per~cent across
the tested margins. A standing payment of \euro{}31--63/kW-yr closes the modelled gap.

(4) Policy ambition sets the emissions outcome. From a 2025 baseline of about
644~MtCO$_2$, 2050 operational emissions reach 352~MtCO$_2$ under Current~Policies,
231 under Stated~Policies, 10 under Net~Zero and 123 under H2~Push. The differences are
driven principally by the fossil-phaseout lever. The heat-pump share climbs from 29 to 49~per~cent across the four, district heat reaches 31~per~cent at about \euro{}96.5/MWh in the deep pathway, and the dispatch screen supports hydrogen at 0, 0.6, 7.5 and 9.0~per~cent of useful heat.

\subsection{Policy implications}

Four implications follow, each scoped to the model's cost criterion. (1) Heat pumps lead by the widest margin across the warm south and east, undercutting hydrogen by about \euro{}99/MWh in Malta and \euro{}88 in Cyprus. (2) Passing ETS2 through to end
users preserves the switching signal, the \euro{}11/MWh per \euro{}50/tCO$_2$ wedge
quantified above, which social protection can leave intact.

(3) Where a hydrogen peaker is wanted, no tested market reaches full capital recovery across the assumed scarcity margins, and the modelled technology-neutral comparison procures unabated gas at the capex assumed here, so a payment conditioned on emissions is what selects hydrogen. Those figures let a capacity-mechanism designer weigh what an emissions condition buys against a technology-neutral auction at the same capital cost. Conditioning costs about
\euro{}172/tCO$_2$ of displaced gas generation in the central case, on top of the carbon
price the gas plant already pays. The technology-neutral
outcome reverses only where hydrogen-ready
machines carry a premium below about 18~per~cent over the gas turbine. At the assumed
gas-turbine emissions intensity, the Article~22(4) threshold does not distinguish the two
technologies; selection therefore requires a more restrictive emissions condition or an
equivalent premium (Section~\ref{sec:results:missing}).

(4) Within the model, the national scope-1 target is not reached through the price signal
alone in every country, so the four markets where the cap binds at 2050, and three more
that bind earlier, need a complementary instrument.

\subsection{Limitations and future work}

Seven limitations and their expected directions are detailed in Supplementary
Section~S2.9. The largest is the weather construction described in
Section~\ref{sec:method}, which conditions the power-sector results. The 278~GW fleet, the
139 full-load hours, the \euro{}357~bn shortfall and the finding that no power-sector
peaker reaches full capital recovery in Table~\ref{tab:arenas_body} rest on it, and the
cold snap's assumed length is a choice,
two days giving 254~GW and ten days 302~GW. An adequacy assessment would run thirty or more recorded weather years and derive the correlation from data, as Kittel \citep{Kittel2026Dunkelflaute} does for European renewable droughts \citep{Staffell2016Wind,Pfenninger2016Solar}, which is beyond a single constructed year. Nonetheless, the recovery shortfall survives every weather variant tested, so that finding does not rest on the construction.

Future work should co-optimise renovation depth with technology choice, build the hourly emissions ledger, resolve within-region heterogeneity, and price the cross-sector capacity value a hydrogen turbine might capture. 
